\section{Descripción de Algoritmos}

\subsection{Data structures}

\textbf{Segment tree} --- \textbf{Complexity:} $O(N \text{ de construcción}, \log N \text{ por operación})$ \\
\textit{Usage:} 1. \texttt{STree(int n)} $\rightarrow$ Reserva el espacio para un arreglo de tamaño $n$. 2. \texttt{upd(int p, int v)} $\rightarrow$ Reemplaza el elemento en la posición \texttt{p} por el valor \texttt{v}. 3. \texttt{query(int l, int r)} $\rightarrow$ Devuelve el resultado (ej. mínimo o suma) de los elementos en el rango $[l, r)$. \\
\textit{Note:} Útil para hacer consultas rápidas sobre un segmento de un arreglo cuando los elementos individuales cambian constantemente. Usa indexación base 0. Para cambiar la operación (ej. de mínimo a máximo), adapta el elemento neutro y la función de combinación.
\vspace{0.15cm}

\textbf{Segment tree lazy} --- \textbf{Complexity:} $O(N \text{ de construcción}, \log N \text{ por actualización/consulta})$ \\
\textit{Usage:} 1. \texttt{STree(int n)} $\rightarrow$ Crea el árbol preparado para "lazy propagation". 2. \texttt{upd(int l, int r, int v)} $\rightarrow$ Aplica la modificación \texttt{v} (ej. sumar \texttt{v}) a todos los elementos dentro del rango $[l, r)$ al mismo tiempo. 3. \texttt{query(int l, int r)} $\rightarrow$ Devuelve el resultado del rango $[l, r)$. \\
\textit{Note:} Igual al Segment Tree, pero este te permite modificar un bloque completo de datos en un solo paso (ej. sumar 5 a todos del índice 2 al 10), sin actualizar celda por celda de manera ineficiente.
\vspace{0.15cm}

\textbf{Persistent segment tree} --- \textbf{Complexity:} $O(\log N \text{ por versión de actualización/consulta})$ \\
\textit{Usage:} 1. \texttt{STree(int n)} $\rightarrow$ Inicializa el árbol. 2. \texttt{upd(int root, int p, int v)} $\rightarrow$ Toma la "foto" del árbol en la versión \texttt{root}, cambia la posición \texttt{p} a \texttt{v}, y devuelve el identificador de la nueva versión creada. 3. \texttt{query(int root, int l, int r)} $\rightarrow$ Consulta el rango $[l, r)$ viajando en el tiempo hacia la versión específica \texttt{root}. \\
\textit{Note:} Es una máquina del tiempo para arreglos. Cada actualización no sobreescribe el pasado, sino que crea una línea temporal alterna. Debes guardar externamente los identificadores devueltos por \texttt{upd} para consultar el pasado.
\vspace{0.15cm}

\textbf{2D segment tree} --- \textbf{Complexity:} $O(NM \text{ de construcción}, \log N \log M \text{ por operación})$ \\
\textit{Usage:} 1. \texttt{build()} $\rightarrow$ Inicializa el árbol leyendo una matriz global predefinida. 2. \texttt{upd(int x, int y, int v)} $\rightarrow$ Cambia el valor en la coordenada $(x, y)$ por \texttt{v}. 3. \texttt{query(int x0, int x1, int y0, int y1)} $\rightarrow$ Devuelve el resultado dentro del rectángulo delimitado por las filas $[x0, x1)$ y columnas $[y0, y1)$. \\
\textit{Note:} Extensión directa para trabajar en grillas 2D. Asegúrate de configurar las variables globales $N$ y $M$ (dimensiones de la matriz) antes de usarlo.
\vspace{0.15cm}

\textbf{Sparse table} --- \textbf{Complexity:} $O(N \log N \text{ de preálculo}, O(1) \text{ por consulta})$ \\
\textit{Usage:} 1. \texttt{st\_init(int *a)} $\rightarrow$ Procesa el arreglo inicial \texttt{a} para preparar las respuestas. 2. \texttt{st\_query(int s, int e)} $\rightarrow$ Devuelve la respuesta instantánea en el rango $[s, e)$. \\
\textit{Note:} Estructura estática. Solo úsala si el arreglo \textbf{nunca} cambia y necesitas responder consultas instantáneamente. Exclusivo para operaciones donde repetir elementos no altera el resultado (idempotentes: min, max, gcd). No sirve para sumas.
\vspace{0.15cm}

\textbf{Fenwick tree} --- \textbf{Complexity:} $O(\log N \text{ por actualización/consulta})$ \\
\textit{Usage:} 1. \texttt{upd(int i0, int v)} $\rightarrow$ Le SUMA el valor \texttt{v} a lo que ya exista en la posición \texttt{i0}. 2. \texttt{get(int i0)} $\rightarrow$ Devuelve la suma acumulada desde el inicio hasta la posición antes de \texttt{i0}. 3. \texttt{get\_sum(int i0, int i1)} $\rightarrow$ Suma del rango $[i0, i1)$. \\
\textit{Note:} Es la versión ligera y rápida de un Segment Tree. Excelente cuando solo necesitas sumas de prefijos y actualizaciones de un solo elemento.
\vspace{0.15cm}

\textbf{Wavelet tree} --- \textbf{Complexity:} $O(N \log \sigma \text{ de construcción}, \log \sigma \text{ por consulta})$ \\
\textit{Usage:} 1. \texttt{WT(int n)} $\rightarrow$ Inicializa para $n$ elementos. 2. \texttt{add(int v)} $\rightarrow$ Agrega un valor al final. 3. \texttt{range(int k, int s, int e, int a, int b)} $\rightarrow$ Cuenta cuántos números caen en el rango de valores $[a, b)$ buscando únicamente dentro del subarreglo desde \texttt{s} hasta \texttt{e}. \\
\textit{Note:} Como un segment tree pero separa valores por bits. Ideal para consultas tipo "cuántos números menores a X hay en este rango del arreglo". Debes comprimir valores muy grandes antes de insertarlos.
\vspace{0.15cm}

\textbf{STL extended set} --- \textbf{Complexity:} $O(\log N \text{ por operación})$ \\
\textit{Usage:} 1. \texttt{ordered\_set} $\rightarrow$ Úsalo como un \texttt{set} de C++ normal. 2. \texttt{find\_by\_order(k)} $\rightarrow$ Devuelve un iterador al \texttt{k}-ésimo elemento más pequeño (0-indexado). 3. \texttt{order\_of\_key(x)} $\rightarrow$ Cuenta exactamente cuántos elementos en el set son estrictamente menores que \texttt{x}. \\
\textit{Note:} Es una caja negra maravillosa que le añade acceso por índice a un Set normal. Requiere directivas GNU de C++, por lo que solo compila en GCC (como en casi todos los jueces).
\vspace{0.15cm}

\textbf{STL rope} --- \textbf{Complexity:} $O(\log N \text{ por separación/inserción/borrado})$ \\
\textit{Usage:} 1. \texttt{rope<T>} $\rightarrow$ Declara tu estructura como si fuera un string o vector. 2. \texttt{insert(pos, value)} $\rightarrow$ Inserta \texttt{value} en la posición \texttt{pos} de manera eficiente. 3. \texttt{erase(pos, len)} $\rightarrow$ Borra \texttt{len} elementos a partir de \texttt{pos}. \\
\textit{Note:} Alternativa para no programar un Treap. Útil si necesitas cortar, pegar y modificar grandes pedazos de un arreglo o cadena de texto. Requiere GNU C++.
\vspace{0.15cm}

\textbf{Treap as BST} --- \textbf{Complexity:} $O(\log N \text{ esperado por operación})$ \\
\textit{Usage:} 1. \texttt{split(pitem t, int key, pitem\& l, pitem\& r)} $\rightarrow$ Parte el árbol en dos: los menores a \texttt{key} en \texttt{l}, y el resto en \texttt{r}. 2. \texttt{insert(pitem\& t, pitem it)} $\rightarrow$ Inserta el nodo \texttt{it}. 3. \texttt{kth(pitem t, int k)} $\rightarrow$ Encuentra el \texttt{k}-ésimo menor en el árbol. \\
\textit{Note:} Es un árbol de búsqueda binaria (como un map/set) balanceado mediante azar. Debes asignar una prioridad aleatoria a cada nodo al crearlo.
\vspace{0.15cm}

\textbf{Implicit treap} --- \textbf{Complexity:} $O(\log N \text{ esperado por operación})$ \\
\textit{Usage:} 1. \texttt{merge(pitem\& t, pitem l, pitem r)} $\rightarrow$ Pega el arreglo \texttt{r} justo a la derecha del arreglo \texttt{l}. 2. \texttt{split(pitem t, pitem\& l, pitem\& r, int sz)} $\rightarrow$ Corta el arreglo dejando los primeros \texttt{sz} elementos en \texttt{l} y el resto en \texttt{r}. 3. \texttt{output(pitem t)} $\rightarrow$ Imprime o recupera el arreglo resultante en orden. \\
\textit{Note:} Funciona como un arreglo mágico en el que puedes cortar un pedazo del medio y pegarlo al inicio en tiempo logarítmico. Las claves son implícitas según la posición.
\vspace{0.15cm}

\textbf{Implicit treap with father} --- \textbf{Complexity:} $O(\log N \text{ esperado por operación})$ \\
\textit{Usage:} Igual al Implicit treap normal, pero incluye punteros al padre. 1. \texttt{root(pitem t, int\& pos)} $\rightarrow$ A partir de un nodo aislado, escala hacia arriba para encontrar la raíz y determina qué posición exacta ocupa este nodo dentro del arreglo completo. \\
\textit{Note:} Útil cuando el problema te da una referencia directa a un elemento y te pregunta "¿en qué posición está actualmente este elemento después de tantas mezclas?".
\vspace{0.15cm}

\textbf{Link-cut tree} --- \textbf{Complexity:} $O(\log N \text{ amortizado por operación})$ \\
\textit{Usage:} 1. \texttt{evert(pitem v)} $\rightarrow$ Gira el árbol para que \texttt{v} sea la nueva raíz suprema. 2. \texttt{link(pitem v, pitem w)} $\rightarrow$ Dibuja una nueva arista conectando \texttt{v} con \texttt{w}. 3. \texttt{cut(pitem v, pitem w)} $\rightarrow$ Destruye la arista entre \texttt{v} y \texttt{w}. \\
\textit{Note:} Mantiene un bosque de árboles que cambia dinámicamente. Útil cuando te piden añadir o borrar aristas sobre la marcha y responder preguntas de conectividad.
\vspace{0.15cm}

\textbf{Link-cut tree with path aggregate} --- \textbf{Complexity:} $O(\log N \text{ amortizado por operación})$ \\
\textit{Usage:} 1. \texttt{mkR(Node x)} $\rightarrow$ Hace a \texttt{x} raíz. 2. \texttt{link(Node x, Node y)} $\rightarrow$ Conecta nodos. 3. \texttt{query(Node x, Node y)} $\rightarrow$ Te da la respuesta combinada (ej. suma, máximo) de todos los nodos en el camino más corto entre \texttt{x} e \texttt{y}. \\
\textit{Note:} Una versión del Link-Cut Tree equipada para responder sobre caminos. Configura las funciones agregadas según lo que necesites calcular.
\vspace{0.15cm}

\textbf{Static convex hull trick} --- \textbf{Complexity:} $O(\log N \text{ por consulta, tras inserción monótona})$ \\
\textit{Usage:} 1. \texttt{CHT::add(Line l)} $\rightarrow$ Agrega una nueva ecuación de recta $y = mx + b$. Ojo: las pendientes $m$ deben insertarse estrictamente ordenadas. 2. \texttt{CHT::query(tc x)} $\rightarrow$ Dado un valor de $x$, evalúa en todas las rectas insertadas y te devuelve el valor de $y$ más pequeño. \\
\textit{Note:} Ideal para optimizar Programación Dinámica. Si buscas máximos en vez de mínimos, multiplica tus pendientes y constantes por $-1$, y multiplica el resultado por $-1$.
\vspace{0.15cm}

\textbf{Dynamic convex hull trick} --- \textbf{Complexity:} $O(\log N \text{ por inserción/consulta})$ \\
\textit{Usage:} 1. \texttt{HullDynamic::add\_line(tc m, tc b)} $\rightarrow$ Agrega una recta en cualquier orden, sin restricciones. 2. \texttt{HullDynamic::query(tc x)} $\rightarrow$ Devuelve el valor \textbf{máximo} al evaluar $x$. \\
\textit{Note:} Funciona igual que el estático pero no requiere que las pendientes lleguen ordenadas. Por defecto, esta implementación busca el MÁXIMO. Para buscar mínimos, inserta $(-m, -b)$ y niega la respuesta.
\vspace{0.15cm}

\textbf{Li Chao tree} --- \textbf{Complexity:} $O(\log C \text{ por inserción/consulta})$ \\
\textit{Usage:} 1. \texttt{LiChaoTree(T xmin, T xmax)} $\rightarrow$ Define el rango válido de valores $x$ que podrías consultar. 2. \texttt{add\_line(Line nw)} $\rightarrow$ Inserta una recta o segmento de recta. 3. \texttt{query(T x)} $\rightarrow$ Encuentra el menor valor de $y$ posible en el punto $x$ evaluando todas las funciones insertadas. \\
\textit{Note:} Es una alternativa moderna y más fácil de modificar que el Convex Hull Trick dinámico.
\vspace{0.15cm}

\textbf{Gain-cost set} --- \textbf{Complexity:} $O(\log N \text{ amortizado por inserción/consulta})$ \\
\textit{Usage:} 1. \texttt{GCS::add(ii x)} $\rightarrow$ Agrega una opción descrita por un par y automáticamente descarta las opciones dominadas (peores en todo). 2. \texttt{GCS::query(int v)} $\rightarrow$ Calcula el mejor rendimiento dado tu recurso \texttt{v}. \\
\textit{Note:} Mantiene una "Frontera de Pareto" limpia usando un set de C++.
\vspace{0.15cm}

\textbf{Disjoint intervals} --- \textbf{Complexity:} $O(K \log N \text{ por inserción})$ \\
\textit{Usage:} 1. \texttt{disjoint\_intervals::add(ii v)} $\rightarrow$ Agrega un nuevo intervalo de tiempo o espacio al sistema. Si choca o se empalma con intervalos ya existentes, los fusiona automáticamente. \\
\textit{Note:} Ideal para colapsar un montón de rangos superpuestos y saber exactamente cuáles son las islas de cobertura resultantes.
\vspace{0.15cm}

\subsection{Graphs}

\textbf{Topological sort} --- \textbf{Complexity:} $O(V+E)$ \\
\textit{Usage:} 1. \texttt{tsort()} $\rightarrow$ Ordena todos los nodos del grafo dirigido de tal forma que ninguna arista apunte hacia atrás. Si hay empates, devuelve el orden lexicográficamente más pequeño. \\
\textit{Note:} Resuelve problemas de dependencia ("hacer A antes que B"). Si el arreglo resultante tiene menos nodos que $V$, significa que hay un ciclo (paradoja).
\vspace{0.15cm}

\textbf{Kruskal MST} --- \textbf{Complexity:} $O(E \log E)$ \\
\textit{Usage:} 1. \texttt{uf\_init()} $\rightarrow$ Prepara los nodos para Union-Find. 2. Ordena tu lista global de aristas de menor a mayor peso. 3. \texttt{kruskal()} $\rightarrow$ Recorre las aristas y devuelve el costo total de conectar todos los nodos sin formar ciclos. \\
\textit{Note:} Árbol de Expansión Mínima. Úsalo para tender redes eléctricas de mínimo costo que conecten todas las ciudades.
\vspace{0.15cm}

\textbf{Dijkstra} --- \textbf{Complexity:} $O((V+E) \log V)$ \\
\textit{Usage:} 1. \texttt{dijkstra(int x)} $\rightarrow$ Explora el mapa buscando la ruta más corta desde $x$ hacia el resto de nodos. El arreglo \texttt{dist[]} guardará los costos mínimos. \\
\textit{Note:} El clásico algoritmo de rutas. Exige que ninguna calle tenga peso negativo. Nodos inaccesibles quedarán marcados con infinito (o -1 según la implementación).
\vspace{0.15cm}

\textbf{Bellman-Ford} --- \textbf{Complexity:} $O(VE)$ \\
\textit{Usage:} 1. \texttt{bford(int src)} $\rightarrow$ Calcula distancias desde \texttt{src}. Propaga adecuadamente las reducciones por ciclos negativos. \\
\textit{Note:} Es como Dijkstra pero soporta costos negativos. Útil para problemas con máquinas del tiempo o intercambios con ganancia neta. Detecta ciclos infinitos de recolección.
\vspace{0.15cm}

\textbf{Floyd-Warshall} --- \textbf{Complexity:} $O(V^3)$ \\
\textit{Usage:} 1. Inicializa tu matriz \texttt{g[u][v]} con el peso directo de cada arista y cero en la diagonal. 2. \texttt{floyd()} $\rightarrow$ Transforma la matriz en las distancias más cortas absolutas entre TODO par de nodos. \\
\textit{Note:} Rápido de programar pero muy lento en tiempo. Úsalo solo para grafos pequeños ($V \le 400$).
\vspace{0.15cm}

\textbf{SCC and 2-SAT} --- \textbf{Complexity:} $O(V+E)$ \\
\textit{Usage:} 1. \texttt{scc()} $\rightarrow$ Agrupa los nodos en Componentes Fuertemente Conexas. 2. \texttt{addor(int a, int b)} $\rightarrow$ Agrega una restricción tipo "o A o B deben cumplirse". 3. \texttt{satisf(int nvar)} $\rightarrow$ Devuelve verdadero si existe una configuración lógica que cumpla todas las reglas. \\
\textit{Note:} La herramienta para problemas de decisiones binarias encadenadas, como rompecabezas de restricciones donde no puedes tener contradicciones lógicas.
\vspace{0.15cm}

\textbf{Articulation, bridges and biconnected} --- \textbf{Complexity:} $O(V+E)$ \\
\textit{Usage:} 1. \texttt{doit()} $\rightarrow$ Analiza el grafo. Al terminar, la bandera \texttt{edge.bridge} indica si destruir esa arista específica partiría la red de comunicación en pedazos aislados. \\
\textit{Note:} Sirve para hallar los puntos débiles de una estructura o una red de servidores.
\vspace{0.15cm}

\textbf{Chu-Liu Edmonds arborescence} --- \textbf{Complexity:} $O(VE)$ \\
\textit{Usage:} 1. \texttt{ChuLiu(int n, int r)} $\rightarrow$ Prepara instancia. 2. \texttt{doit()} $\rightarrow$ Devuelve el costo de la red mínima dirigida desde la raíz $r$. \\
\textit{Note:} Es el análogo de Kruskal (MST) pero para calles de UN SOLO sentido. Garantiza rutas válidas dirigidas desde la base $r$ hacia todo el mundo.
\vspace{0.15cm}

\textbf{LCA binary lifting} --- \textbf{Complexity:} $O(N \log N \text{ preálculo}, \log N \text{ por consulta})$ \\
\textit{Usage:} 1. \texttt{lca\_init()} $\rightarrow$ Recorre el árbol y anota saltos de tamaño potencia de 2. 2. \texttt{lca(int x, int y)} $\rightarrow$ Devuelve el nodo Ancestral Común más profundo entre $x$ e $y$. \\
\textit{Note:} Vital para resolver distancias reales en árboles usando la fórmula: $dist(u,v) = \text{prof}(u) + \text{prof}(v) - 2 \cdot \text{prof}(LCA(u,v))$.
\vspace{0.15cm}

\textbf{Heavy-light decomposition} --- \textbf{Complexity:} $O(N \text{ inicialización}, \log^2 N \text{ por consulta})$ \\
\textit{Usage:} 1. \texttt{hld\_init()} $\rightarrow$ Acomoda el árbol en un arreglo recto. 2. \texttt{query(int x, int y, STree\& rmq)} $\rightarrow$ Consulta el segmento continuo de las ramas del árbol apoyándote en un Segment Tree. \\
\textit{Note:} Herramienta definitiva para aplicar modificaciones o cálculos sobre trayectorias dentro de un árbol como si fueran segmentos de un arreglo 1D.
\vspace{0.15cm}

\textbf{Centroid decomposition} --- \textbf{Complexity:} $O(N \log N)$ \\
\textit{Usage:} 1. \texttt{centroid()} $\rightarrow$ Encuentra el "centro de gravedad" del árbol y lo extrae recursivamente. 2. Dentro de \texttt{cdfs}, programas lo que deseas calcular con las mitades particionadas. \\
\textit{Note:} Divide el problema de conteo de caminos en mitades. Usado para problemas tipo "¿cuántos caminos en el árbol suman $K$ de longitud?".
\vspace{0.15cm}

\textbf{Parallel DFS} --- \textbf{Complexity:} $O(V+E)$ \\
\textit{Usage:} 1. \texttt{Tree::go(int x)} $\rightarrow$ Ejecuta una exploración simultánea. \\
\textit{Note:} Permite recorrer y comparar varios árboles/grafos a la vez manejando su estado conjunto.
\vspace{0.15cm}

\textbf{Eulerian path} --- \textbf{Complexity:} $O(V+E)$ \\
\textit{Usage:} 1. \texttt{add\_edge(int a, int b)} $\rightarrow$ Agrega los trazos. 2. \texttt{get\_path(int x)} $\rightarrow$ Devuelve un arreglo con el orden de vértices a pisar para recorrer TODAS las aristas exactamente una vez. \\
\textit{Note:} El problema de los "Puentes de Königsberg" o "dibujar sin levantar el lápiz". Revisa la factibilidad de los grados (par/impar) de los nodos antes de llamarla.
\vspace{0.15cm}

\textbf{Offline dynamic connectivity} --- \textbf{Complexity:} $O((N+Q) \log Q)$ \\
\textit{Usage:} 1. \texttt{add(int l, int r, Query e)} $\rightarrow$ Especifica entre qué turnos $[l, r)$ existe cierta arista. 2. \texttt{go(int k, int s, int e)} $\rightarrow$ Ejecuta el historial temporal usando rollback sobre un DSU. \\
\textit{Note:} Caja negra para procesar queries de grafos que sufren alteraciones y borrados. Requiere guardar todas las consultas de antemano ("offline").
\vspace{0.15cm}

\textbf{Dominator tree} --- \textbf{Complexity:} $O((V+E) \alpha(V))$ \\
\textit{Usage:} 1. \texttt{dominators(int rt)} $\rightarrow$ Rellena un arreglo \texttt{idom[v]} que indica qué vértice actúa como aduana obligatoria y directa para llegar a $v$ desde $rt$. \\
\textit{Note:} Halla nodos críticos inescapables en grafos dirigidos.
\vspace{0.15cm}

\textbf{Edmonds blossom} --- \textbf{Complexity:} $O(V^2 E)$ \\
\textit{Usage:} 1. \texttt{edmonds()} $\rightarrow$ Devuelve el tamaño máximo de emparejamientos sin superposición que se pueden formar en el grafo. \\
\textit{Note:} Emparejamiento máximo para "Grafos Generales". Úsalo cuando las reglas te impiden dividir los nodos en dos grupos (bipartito) separados.
\vspace{0.15cm}

\subsection{Math}

\textbf{Integer floor division} --- \textbf{Complexity:} $O(1)$ \\
\textit{Usage:} 1. \texttt{floordiv(ll x, ll y, ll\& q, ll\& r)} $\rightarrow$ Calcula la división hacia $-\infty$ y deposita el cociente en $q$ y el residuo positivo/concordante en $r$. \\
\textit{Note:} C++ estándar trunca hacia $0$, lo que genera errores garrafales con coordenadas negativas. Usa esto para que las matemáticas modulares cuadren perfecto sin importar el signo.
\vspace{0.15cm}

\textbf{Extended Euclid} --- \textbf{Complexity:} $O(\log \min(a,b))$ \\
\textit{Usage:} 1. \texttt{euclid(ll a, ll b, ll\& x, ll\& y)} $\rightarrow$ Calcula el Máximo Común Divisor (GCD) entre $a$ y $b$, y averigua los factores $x, y$ para resolver $ax + by = \text{GCD}$. \\
\textit{Note:} Imprescindible en la búsqueda del Inverso Modular y Teorema Chino del Resto.
\vspace{0.15cm}

\textbf{Sieve of Eratosthenes} --- \textbf{Complexity:} $O(N \log \log N)$ \\
\textit{Usage:} 1. \texttt{init\_sieve()} $\rightarrow$ Elimina múltiplos velozmente dejando registrados únicamente los números primos hasta un límite máximo (\texttt{MAXN}). \\
\textit{Note:} Llama esta función una sola vez al inicio. Ideal cuando tienes múltiples casos de prueba que requieren factorizar o preguntar la primalidad de números pequeños-medianos ($<10^7$).
\vspace{0.15cm}

\textbf{Generate divisors} --- \textbf{Complexity:} $O(d(n))$ \\
\textit{Usage:} 1. \texttt{divisors(vector<pair<ll,int>> f)} $\rightarrow$ A partir de la huella prima (ej. $2^2 \times 3^1$), genera la lista completa de todos los divisores enteros posibles. \\
\textit{Note:} Útil y rápida cuando necesitas saber entre cuántas partes exactas se puede dividir un número ya factorizado.
\vspace{0.15cm}

\textbf{Pollard Rho factorization} --- \textbf{Complexity:} $O(N^{1/4})$ \\
\textit{Usage:} 1. \texttt{rabin(ll n)} $\rightarrow$ Test de primalidad para números ridículamente grandes ($10^{18}$). 2. \texttt{fact(ll n, map<ll,int>\& f)} $\rightarrow$ Desmenuza un número gigante guardando su factorizacion en el mapa $f$. \\
\textit{Note:} La salvación absoluta cuando la criba de Eratóstenes explota por falta de RAM. Usa aleatoriedad para hallar los divisores más escurridizos al instante.
\vspace{0.15cm}

\textbf{Simpson integration} --- \textbf{Complexity:} $O(N)$ \\
\textit{Usage:} 1. \texttt{integrate(double f(double), double a, double b, int n)} $\rightarrow$ Aproxima el área bajo la curva de la función matemática $f$ usando $n$ subintervalos. \\
\textit{Note:} Usa $n$ par y grande ($100000$) para excelente precisión en problemas de áreas continuas y curvas de cálculo.
\vspace{0.15cm}

\textbf{Polynomial utilities} --- \textbf{Complexity:} $O(\text{Depende})$ \\
\textit{Usage:} 1. \texttt{poly<T>} $\rightarrow$ Úsalo como tipo de dato de tus ecuaciones. 2. \texttt{roots(poly<> p)} $\rightarrow$ Halla raíces enteras. 3. \texttt{interpolate(x, y)} $\rightarrow$ Halla la ecuación que cruza dichos puntos. \\
\textit{Note:} Permite manipular polinomios sin complicarte con índices manuales.
\vspace{0.15cm}

\textbf{Bairstow roots} --- \textbf{Complexity:} $O(\text{ITER} \cdot \text{Grado})$ \\
\textit{Usage:} 1. \texttt{bairstow(poly<> p)} $\rightarrow$ Retira y calcula aproximaciones cuadráticas del polinomio mayor. \\
\textit{Note:} Método numérico iterativo potente para descifrar raíces complejas o reales reduciendo el polinomio.
\vspace{0.15cm}

\textbf{Fast Hadamard Transform} --- \textbf{Complexity:} $O(N \log N)$ \\
\textit{Usage:} 1. \texttt{fht(ll* p, int n, bool inv)} $\rightarrow$ Transforma el arreglo al dominio de Walsh. 2. \texttt{multiply(vector<ll>\& p1, vector<ll>\& p2)} $\rightarrow$ Convoluciona los elementos operando bajo reglas bit a bit (XOR, AND, OR). \\
\textit{Note:} Sirve para responder rápidamente combinaciones de bits (ej. "de cuántas formas saco el número C sumando un ítem A y un ítem B usando la operación XOR"). Modifica el símbolo interno según la operación deseada.
\vspace{0.15cm}

\textbf{Karatsuba multiplication} --- \textbf{Complexity:} $O(N^{1.585})$ \\
\textit{Usage:} 1. \texttt{multiply(vector<tp> p0, vector<tp> p1)} $\rightarrow$ Multiplica (convoluciona) dos polinomios/números masivos rompiéndolos recursivamente. \\
\textit{Note:} La alternativa al FFT. Úsalo si tu entorno tiene problemas con módulos y números flotantes que arruinan la precisión del FFT tradicional.
\vspace{0.15cm}

\textbf{Linear Diophantine equations} --- \textbf{Complexity:} $O(\log \min(a,b))$ \\
\textit{Usage:} 1. \texttt{diophantine(ll a, ll b, ll r)} $\rightarrow$ Encuentra los valores $X$ e $Y$ enteros que satisfacen $A \cdot X + B \cdot Y = R$. \\
\textit{Note:} La condición de oro es que $R$ sea obligatoriamente un múltiplo divisible por el $\text{GCD}(A,B)$, en caso contrario no existirá solución entera.
\vspace{0.15cm}

\textbf{Modular inverse} --- \textbf{Complexity:} $O(\log \text{mod})$ \\
\textit{Usage:} 1. \texttt{inv(ll a, ll mod)} $\rightarrow$ Encuentra el número por el que debes multiplicar para simular una división entre $a$ bajo la aritmética de un módulo de reloj. \\
\textit{Note:} $a$ y el \texttt{mod} deben ser coprimos ($\text{GCD}=1$) para que el inverso exista.
\vspace{0.15cm}

\textbf{Simplex Algorithm} --- \textbf{Complexity:} $\text{Exponencial en teoría, muy rápido en práctica}$ \\
\textit{Usage:} 1. Traduce tu problema a ecuaciones matemáticas, cárgalas en la matriz de la función y recibe el máximo (o mínimo) valor posible de tu función objetivo. \\
\textit{Note:} La bala de plata para "Tengo X recursos y condiciones variadas; quiero maximizar mi ganancia monetaria sin violar los límites". (Programación Lineal).
\vspace{0.15cm}

\textbf{Discrete logarithm} --- \textbf{Complexity:} $O(\sqrt{M})$ \\
\textit{Usage:} 1. \texttt{discrete\_log(ll a, ll b, ll m)} $\rightarrow$ Resuelve el logaritmo en un reloj: encuentra cuántas veces se debe potenciar $a$ para que resulte ser $b$ bajo módulo $m$. \\
\textit{Note:} Utiliza el ingenioso algoritmo Baby-Step Giant-Step para acortar el tiempo de $M$ a su raíz cuadrada.
\vspace{0.15cm}

\textbf{Berlekamp-Massey} --- \textbf{Complexity:} $O(N^2)$ \\
\textit{Usage:} 1. \texttt{BM(vi x)} $\rightarrow$ Dale un puñado de números de una secuencia repetitiva desconocida, y este algoritmo deducirá la relación/fórmula lineal que la produce y sus coeficientes exactos. \\
\textit{Note:} Adulteración de DP puro. Si no logras pensar la lógica del problema, usa backtracking bruto, escupe los primeros términos de respuesta, mételos en BM y te dará el DP de regalo.
\vspace{0.15cm}

\textbf{Linear recurrence} --- \textbf{Complexity:} $O(K^2 \log N)$ \\
\textit{Usage:} 1. \texttt{LinearRec} y \texttt{calc(ll k)} $\rightarrow$ Dado tu estado inicial y los coeficientes de recurrencia (obtenidos a mano o por BM), halla instantáneamente el término súper-lejos (ej. $10^{18}$) de la sucesión. \\
\textit{Note:} Trabaja bajo la misma teoría mágica que Exponenciación de Matrices pero de forma polinómica.
\vspace{0.15cm}

\textbf{Tonelli-Shanks} --- \textbf{Complexity:} $O(\log^2 P)$ \\
\textit{Usage:} 1. \texttt{tonelli\_shanks(ll n, ll p)} $\rightarrow$ Dado un número $n$ y un módulo primo $p$, extrae el valor $X$ que cumple $X^2 \equiv n \pmod p$. \\
\textit{Note:} Tu calculadora para Raíz Cuadrada en aritmética modular. Si la función \texttt{legendre} devuelve -1, significa que no existe ninguna raíz bajo las reglas de ese módulo.
\vspace{0.15cm}

\textbf{Prime counting} --- \textbf{Complexity:} $O(N^{2/3} / \log^{1/3} N)$ \\
\textit{Usage:} 1. \texttt{count\_primes(ll n)} $\rightarrow$ Devuelve exactamente el número de primos menores o iguales a $n$. \\
\textit{Note:} Espectacular algoritmo para cuando $n$ es exagerado ($10^{11}$) y Eratóstenes falla, pero el problema es simple: "dime la cantidad de primos".
\vspace{0.15cm}

\subsection{Geometry}

\textbf{Point / Vector operations} --- \textbf{Complexity:} $O(1)$ \\
\textit{Usage:} 1. \texttt{cross(p1, p2)} $\rightarrow$ Producto Cruz; calcula áreas y dictamina si el polígono gira a izquierda o derecha. 2. \texttt{dot(p1, p2)} $\rightarrow$ Producto Punto; calcula proyecciones. \\
\textit{Note:} La base de la vida geométrica. El producto cruz salva la vida al esquivar funciones de ángulos con coma flotante, evitando errores catastróficos de aproximación.
\vspace{0.15cm}

\textbf{Line / Segment operations} --- \textbf{Complexity:} $O(1)$ \\
\textit{Usage:} 1. \texttt{operator\string^(ln l, ln m)} $\rightarrow$ Ejecuta y extrae exactamente el punto de intersección entre las dos rectas. \\
\textit{Note:} Mantén un margen \texttt{EPS} sano para no comparar \texttt{0.0} exacto por temas de decimales flotantes.
\vspace{0.15cm}

\textbf{Circle} --- \textbf{Complexity:} $O(1 \text{ o } N \text{ para interacciones})$ \\
\textit{Usage:} 1. \texttt{circle(pt o, double r)} $\rightarrow$ Crea y clasifica círculos. 2. \texttt{intercircles(vector<circle> c)} $\rightarrow$ Estima el área consolidada de múltiples círculos superpuestos como en un diagrama de Venn masivo. \\
\textit{Note:} Altamente sensible al parámetro \texttt{EPS}.
\vspace{0.15cm}

\textbf{Polygon} --- \textbf{Complexity:} $O(N)$ \\
\textit{Usage:} 1. \texttt{pol::area()} $\rightarrow$ Calcula el área orientada internamente dividiendo todo en triángulos desde el origen mediante el producto cruz. 2. \texttt{pol::has(pt q)} $\rightarrow$ Confirma si el punto $q$ se encuentra dentro de los dominios del polígono. \\
\textit{Note:} Revisa que la lista de vértices esté ordenada y respete la convexidad y dirección que espera tu código.
\vspace{0.15cm}

\textbf{Plane} --- \textbf{Complexity:} $O(1)$ \\
\textit{Usage:} 1. \texttt{plane::side(pt q)} $\rightarrow$ Verifica en qué cara o hemisferio respecto al plano infinito se halla el punto $q$. \\
\textit{Note:} Para aquellos infames problemas que obligan a saltar al espacio 3D.
\vspace{0.15cm}

\textbf{Radial order} --- \textbf{Complexity:} $O(N \log N)$ \\
\textit{Usage:} 1. Aplícalo como un comparador estricto para C++ \texttt{std::sort} y ordenará los puntos geométricos alrededor del centro barriendo angularmente. \\
\textit{Note:} Útil y necesario para construir Envolventes Convexos (Convex Hulls) y técnicas de "Sweep Line Angular" sin usar el impreciso \texttt{atan2}.
\vspace{0.15cm}

\textbf{Minkowski sum} --- \textbf{Complexity:} $O(N+M)$ \\
\textit{Usage:} 1. \texttt{minkowski\_sum(vector<pt> p, vector<pt> q)} $\rightarrow$ Agrupa dos figuras sólidas en una sola silueta que representa todos los trazos vectores posibles. \\
\textit{Note:} La varita mágica para detectar colisiones o choques exactos entre dos polígonos convexos trasladándose uno hacia el otro.
\vspace{0.15cm}

\textbf{Delaunay triangulation} --- \textbf{Complexity:} $O(N \log N \text{ esperado})$ \\
\textit{Usage:} 1. \texttt{delaunay(vector<pt> v)} $\rightarrow$ Ensambla triángulos sanos, evitando ángulos puntiagudos inútiles, para construir una red entre puntos dispersos. \\
\textit{Note:} Es la estructura dual al Diagrama de Voronoi. Evita mandarle puntos encimados o idénticos.
\vspace{0.15cm}

\subsection{Strings}

\textbf{KMP} --- \textbf{Complexity:} $O(N+M)$ \\
\textit{Usage:} 1. \texttt{kmppre(string\& t)} $\rightarrow$ Estudia la meta. 2. \texttt{kmp(string\& s, string\& t)} $\rightarrow$ Busca rápidamente sin repetir pasajes inútiles la existencia del patrón $t$ a lo largo de todo $s$. \\
\textit{Note:} La caja original del String Matching ("Ctrl+F"). Excelente velocidad y bajo consumo de memoria.
\vspace{0.15cm}

\textbf{Z function} --- \textbf{Complexity:} $O(N)$ \\
\textit{Usage:} 1. \texttt{z\_function(string\& s)} $\rightarrow$ Genera el arreglo $Z$. Cada celda \texttt{z[i]} indica cuántos caracteres del texto a partir de allí son réplicas exactas del prefijo mismo de toda la cadena. \\
\textit{Note:} Sumamente poderoso para buscar patrones si unes el patrón al inicio del texto usando un carácter especial ciego (como `$\#$`) en el medio.
\vspace{0.15cm}

\textbf{Manacher} --- \textbf{Complexity:} $O(N)$ \\
\textit{Usage:} 1. \texttt{manacher(string\& s)} $\rightarrow$ Analiza la cadena insertando delimitadores invisibles para escupir instantáneamente todos los palíndromos alojados sin importar su paridad. \\
\textit{Note:} Extrae el palíndromo más largo o cuenta los subpalíndromos centrados en cualquier carácter velozmente.
\vspace{0.15cm}

\textbf{Suffix array} --- \textbf{Complexity:} $O(N \log N)$ \\
\textit{Usage:} 1. \texttt{constructSA(string\& s)} $\rightarrow$ Retorna el índice catalogado, ordenado de la A a la Z, de todas las formas posibles de terminar o diseccionar el texto. \\
\textit{Note:} Sirve como tu guía telefónica definitiva. Al estar ordenado, puedes usar Búsqueda Binaria sobre el arreglo resultante para buscar palabras parciales velozmente.
\vspace{0.15cm}

\textbf{LCP array} --- \textbf{Complexity:} $O(N)$ \\
\textit{Usage:} 1. \texttt{computeLCP(string\& s, vector<int>\& sa)} $\rightarrow$ Genera un arreglo de similitudes entre vecinos en tu Suffix Array. \\
\textit{Note:} Siempre se usa de la mano con Suffix Array. Convierte los misteriosos problemas de texto en operaciones matemáticas sobre arreglos numéricos (ej. hallar el substring común más largo).
\vspace{0.15cm}

\textbf{Suffix tree} --- \textbf{Complexity:} $O(N)$ \\
\textit{Usage:} 1. \texttt{SuffixTree(string s)} $\rightarrow$ Condensa el abrumador universo de combinaciones del texto en un bello árbol estructurado. \\
\textit{Note:} Extremadamente pesado. No olvides rematar el texto con un sello de fin (ej. `\$`) para materializar todas las hojas posibles sin problemas invisibles.
\vspace{0.15cm}

\textbf{String hashing} --- \textbf{Complexity:} $O(N \text{ preálculo}, O(1) \text{ consulta})$ \\
\textit{Usage:} 1. \texttt{Hash(string s)} $\rightarrow$ Fija números de base. 2. \texttt{get(int l, int r)} $\rightarrow$ Otorga el código de identidad inalterable de un pedazo del texto. \\
\textit{Note:} Transforma comparaciones molestas de caracteres letra por letra, en triviales comprobaciones "$\text{hash}_A == \text{hash}_B$" a velocidad de luz.
\vspace{0.15cm}

\subsection{Other}

\textbf{Min-cost max-flow (MCMF)} --- \textbf{Complexity:} $O(\text{Flujo Max} \cdot E \log V)$ \\
\textit{Usage:} 1. \texttt{add\_edge(int u, int v, int cap, int cost)} $\rightarrow$ Conecta un tubo entre \texttt{u} y \texttt{v} donde pueden pasar hasta \texttt{cap} elementos en total, cobrando un peaje o castigo de \texttt{cost} por CADA unidad enviada. 2. \texttt{get\_flow(int s, int t)} $\rightarrow$ Inicia la simulación en \texttt{s} apuntando a \texttt{t}. Devuelve el par consolidado de \texttt{\{Máximo\_Flujo, Mínimo\_Costo\}}. \\
\textit{Note:} Envía toda el agua posible pero busca ávidamente las rutas más económicas primero para ahorrar el total. Si tienes aristas iniciales negativas, deberás incluir la limpieza mágica de los "Potenciales" antes del Dijkstra subyacente.
\vspace{0.15cm}

\textbf{Mo algorithm} --- \textbf{Complexity:} $O((N+Q) \sqrt{N})$ \\
\textit{Usage:} 1. \texttt{mos()} $\rightarrow$ Revuelve y ordena las preguntas (consultas) inteligentemente basadas en bloques de raíz cuadrada para anular los saltos extremos largos en el arreglo. 2. \texttt{add(pos)} y \texttt{remove(pos)} $\rightarrow$ Al mover la cortina (tu lupa de visualización) un pasito, informas qué pasa con la estadística o el contador total. \\
\textit{Note:} Estrictamente "Offline". Útil cuando no puedes usar un Segment Tree porque combinar las respuestas de dos mitades separadas es infernal, pero añadir/quitar un ítem de tu cuenta actual te toma solo $O(1)$.
\vspace{0.15cm}

\textbf{Divide and conquer DP optimization} --- \textbf{Complexity:} $O(K \cdot N \log N)$ \\
\textit{Usage:} 1. \texttt{doit(int s, int e, int s0, int e0, int i)} $\rightarrow$ Corta a la mitad la zona incierta sabiendo el punto de corte óptimo validado. 2. \texttt{doall()} $\rightarrow$ Orquesta los bucles sobre las iteraciones. \\
\textit{Note:} Resuelve misterios de Programación Dinámica tipo "Partición de arreglo en $K$ componentes", acuchillando un inaceptable $O(K \cdot N^2)$ para convertirlo en algo muy veloz. Condición obligatoria: tu punto de corte debe ser estrictamente monótono.
\vspace{0.15cm}

\textbf{Date conversion} --- \textbf{Complexity:} $O(1)$ \\
\textit{Usage:} 1. \texttt{dateToInt(int y, int m, int d)} $\rightarrow$ Exprime una fastidiosa fecha del calendario gregoriano en un simple número absoluto (Día Juliano). \\
\textit{Note:} ¡Huye de programar meses bisiestos y validaciones a mano! Convierte fechas a enteros, réstalos libremente, y aplícales módulo $7$ para develar el día semanal.
\vspace{0.15cm}
